\فصل{نتایج تجربی}

در ادامه‌ی معرفی روش پیشنهادی در فصل قبل، در این فصل نتایج حاصل از ارزیابی تجربی ارائه می‌شود. هدف اصلی این ارزیابی، بررسی عملکرد و پایداری روش در حالتهای متنوعی است که یک مدل در شرایط واقعی با آن مواجه می‌شود. از آنجا که مسئله‌ی اصلی پژوهش حاضر، انطباق مدل در زمان آزمون به صورت برخط و مستمر است، آزمایش‌ها به‌گونه‌ای طراحی شده‌اند که توانایی روش در حفظ دقت، پایداری و سازگاری در مواجهه با داده‌های جدید و تغییرات توزیعی را به شکلی جامع نشان دهند.  

برای دستیابی به این هدف، دو حالتی اصلی مدنظر قرار گرفته است. در حالت اول، مدل به صورت مستقل و بدون دسترسی به داده‌های آموزشی قبلی، صرفاً با استفاده از داده‌های جاری در آزمون سازگار می‌شود. در این حالت، داده‌های آزمون به‌صورت پیوسته اما تنها از یک توزیع مشخص به‌صورت $D_{\text{test}} \sim D_1$ وارد می‌شوند.

در حالت دوم، جریان پیوسته‌ای از داده‌ها در طول زمان در اختیار مدل قرار دارد و توزیع داده‌ها به صورت $D_{\text{test}} \sim D_1, \ldots, D_k$ در مراحل مختلف تغییر می‌کند. در این حالت، مدل باید به‌طور مداوم خود را با شرایط جدید وفق دهد.\\
بررسی عملکرد روش در هر دو حالت، امکان ارزیابی جامع‌تری را فراهم می‌کند و نقاط قوت و ضعف آن را در شرایط متفاوت نمایان می‌سازد.  

یکی از مسائل کلیدی در روش‌های انطباق، حساسیت به فراپارامترها است. انتخاب نامناسب نرخ یادگیری، تعداد به‌روزرسانی‌ها یا سایر پارامترهای تنظیمی می‌تواند موجب کاهش شدید کارایی مدل شود. بنابراین در این فصل، عملکرد روش پیشنهادی در مجموعه‌ای از تنظیمات مختلف مورد آزمایش قرار می‌گیرد تا میزان پایداری و استحکام روش در برابر تغییر فراپارامترها مشخص شود. این بررسی نه تنها اعتبار نتایج را افزایش می‌دهد بلکه نشان می‌دهد روش تا چه اندازه در حالتهای عملی و شرایطی که تنظیم دقیق فراپارامترها دشوار است، قابل اعتماد خواهد بود. اهمیت این موضوع به‌ویژه در حالتی انطباق زمان آزمون برخط بیشتر است، چراکه در این شرایط هیچ دسترسی به داده‌های منبع یا برچسب نمونه‌ها برای تنظیم یا اعتبارسنجی فراپارامترها در زمان آزمون وجود ندارد. در حالی‌که بسیاری از پژوهش‌های پیشین نتایج خود را با تکیه بر تنظیم دقیق فراپارامترها گزارش کرده‌اند، در محیط‌های واقعی برخط چنین امکانی عملاً فراهم نیست. از این رو، بررسی حساسیت روش پیشنهادی به فراپارامترها نقش تعیین‌کننده‌ای در اعتبارسنجی عملی آن دارد.  

از سوی دیگر، در انطباق برخط زمان آزمون مسئله‌ی انباشت خطا و فروپاشی مدل مطرح می‌شود. به‌روزرسانی‌های پیاپی در طول زمان اگر به‌درستی کنترل نشوند، می‌توانند باعث انحراف تدریجی مدل از ساختار اصلی و در نهایت از دست رفتن کارایی شوند. بنابراین، در این فصل علاوه بر بررسی دقت و بهبود عملکرد، توجه ویژه‌ای به رفتار روش در برابر انباشت خطا، پایداری بلندمدت و جلوگیری از فروپاشی معطوف شده است. این مسئله از آن جهت اهمیت دارد که در کاربردهای واقعی، مدل‌ها معمولاً در محیط‌هایی پویا و با داده‌هایی غیرایستا به کار گرفته می‌شوند و دوام عملکرد آن‌ها نقش تعیین‌کننده‌ای در قابلیت استفاده عملی خواهد داشت.  

به طور کلی، ساختار این فصل بر اساس پرسش‌های کلیدی زیر سامان‌دهی شده است:
\begin{itemize}
	\item روش پیشنهادی در حالتهای انطباق زمان آزمون و انطباق مستمر چه عملکردی دارد؟
	\item میزان حساسیت روش به تغییرات فراپارامترها چگونه است و آیا می‌توان آن را روشی پایدار و مقاوم در نظر گرفت؟
	\item در شرایط طولانی‌مدت و تحت به‌روزرسانی‌های مکرر، رفتار مدل از منظر انباشت خطا و پدیده فروپاشی چگونه خواهد بود؟
\end{itemize}

برای پاسخ به این پرسش‌ها، در ادامه ابتدا مجموعه داده‌های مورد استفاده معرفی می‌شوند. سپس معیارهای ارزیابی و جزئیات پیاده‌سازی روش تشریح شده و در نهایت نتایج به‌دست‌آمده همراه با تحلیل‌های مقایسه‌ای ارائه خواهد شد. این فصل با جمع‌بندی نکات اصلی به پایان می‌رسد تا چارچوبی روشن برای ارزیابی کلی روش در اختیار قرار گیرد.

\قسمت{مجموعه داده}

برای ارزیابی روش پیشنهادی، آزمایش‌ها بر روی مجموعه‌داده‌های {\lr{CIFAR-10}، \lr{CIFAR-100} }\cite{cifar10-100} و نسخه‌های مخدوش آن‌ها یعنی {\lr{CIFAR-10-C} و \lr{CIFAR-100-C} }\cite{cifar10c-100c} انجام شده است. 

مجموعه‌داده‌ی \lr{CIFAR-10} شامل ۶۰۰۰۰ تصویر رنگی با ابعاد $32\times32\times3$ است که در ۱۰ کلاس دسته‌بندی شده‌اند (هر کلاس ۶۰۰۰ تصویر). این مجموعه به ۵۰۰۰۰ تصویر برای آموزش و ۱۰۰۰۰ تصویر برای آزمون تقسیم شده است. مجموعه‌داده‌ی \lr{CIFAR-100} ساختاری مشابه دارد، اما شامل ۱۰۰ کلاس است که هر کلاس دارای ۶۰۰ تصویر بوده و بنابراین طبقه‌بندی آن چالش‌برانگیزتر است. 

برای بررسی پایداری مدل در شرایط توزیع‌های متفاوت، نسخه‌های مخدوش این داده‌ها مورد استفاده قرار گرفته‌اند. مجموعه‌های \lr{CIFAR-10-C} و \lr{CIFAR-100-C} از تصاویر آزمون اصلی مشتق شده‌اند که تحت ۱۵ نوع خرابی مختلف (مانند نویز، {بلور‌\پانویس{Blur}}، تغییرات آب‌وهوایی و فشرده‌سازی) اعمال شده‌اند. هر نوع خرابی در پنج سطح شدت \lr{Severity Levels 1–5 } موجود است که مجموعاً ۷۵ حالت مختلف از داده‌های مخدوش را تشکیل می‌دهد. هر یک از این حالت‌ها شامل ۱۰۰۰۰ تصویر است و در نتیجه، کل نسخه‌ی مخدوش برای هر مجموعه شامل ۷۵۰۰۰۰ تصویر می‌شود. این تنوع امکان بررسی دقیق عملکرد روش پیشنهادی در برابر تغییرات شدید توزیع داده‌ها و سنجش میزان استحکام آن را فراهم می‌آورد.  

خلاصه‌ای از ویژگی‌های اصلی این مجموعه‌داده‌ها در جدول \ref{table:datasets} ارائه شده است.  

% \begin{table}[t]
% 	\centering
% 	\caption{خلاصه‌ای از مجموعه‌داده‌های اصلی و مخدوش مورد استفاده در آزمایش‌ها}
% 	\label{table:datasets}
% 	\begin{latin}
% 		\begin{tabular}{lccccc}
% 			\hline
% 			\textbf{Dataset} & \textbf{\#Train} & \textbf{\#Test} & \textbf{\#Corr.} & \textbf{\#Severity} & \textbf{\#Classes} \\
% 			\hline
% 			CIFAR-10      & 50,000 & 10,000 & 1  & 1 & 10  \\
% 			CIFAR-100     & 50,000 & 10,000 & 1  & 1 & 100 \\
% 			CIFAR-10-C    &   --   & 950,000 & 15 & 5 & 10  \\
% 			CIFAR-100-C   &   --   & 950,000 & 15 & 5 & 100 \\
% 			\hline
% 		\end{tabular}
% 	\end{latin}
% \end{table}
\begin{table}[t]
    \centering
    \small
    \caption{\RL{خلاصه‌ای از مجموعه‌داده‌های اصلی و مخدوش مورد استفاده در آزمایش‌ها}}
    \label{table:datasets}
    \setlength{\tabcolsep}{4pt}
    \begin{tabular}{p{4cm} p{2cm} p{2cm} p{2cm} p{2cm} p{2cm}}
    \toprule
    \textbf{\RL{مجموعه داده}} & 
    \textbf{\RL{\#آموزش}} & 
    \textbf{\RL{\#آزمون}} & 
    \textbf{\RL{\#تخریب}} & 
    \textbf{\RL{\#شدت}} & 
    \textbf{\RL{\#کلاس‌ها}} \\ \midrule
    \lr{CIFAR-10}      & \lr{50,000} & \lr{10,000} & \lr{1}  & \lr{1} & \lr{10}  \\
    \lr{CIFAR-100}     & \lr{50,000} & \lr{10,000} & \lr{1}  & \lr{1} & \lr{100} \\
    \lr{CIFAR-10-C}    & \lr{--}     & \lr{950,000} & \lr{15} & \lr{5} & \lr{10}  \\
    \lr{CIFAR-100-C}   & \lr{--}     & \lr{950,000} & \lr{15} & \lr{5} & \lr{100} \\
    \bottomrule
    \end{tabular}
\end{table}



شکل \ref{fig:dataser_severity_levels} نشان‌دهنده‌ی نمونه‌هایی از یک نوع خرابی مشخص (نویز گاوسی) در سطوح شدت مختلف است. همان‌طور که مشاهده می‌شود، با افزایش سطح شدت از ۱ تا ۵، میزان تغییرات و اختلال در تصویر افزایش یافته و چالش مدل برای حفظ دقت طبقه‌بندی بیشتر می‌شود.

\begin{figure}[t]
	\centering
	\subfigure[بدون تخریب]{	
		\includegraphics[width=0.23\linewidth]{images/severity/cifar100c/original_17.png}
	}
	\subfigure[سطح ۱]{	
		\includegraphics[width=0.23\linewidth]{images/severity/cifar100c/severity_1/gaussian_noise_17.png}
	}
	\subfigure[سطح ۳]{	
		\includegraphics[width=0.23\linewidth]{images/severity/cifar100c/severity_3/gaussian_noise_17.png}
	}
	\subfigure[سطح ۵]{
		\includegraphics[width=0.23\linewidth]{images/severity/cifar100c/severity_5/gaussian_noise_17.png}
	}
	\caption[مقایسه‌ی سطوح شدت تخریب نویز گاوسی تصویر در مجموعه داده‌ی  \lr{Cifar100-C}]{مقایسه‌ی سطوح شدت تخریب نویز گاوسی تصویر در مجموعه داده‌ی  \lr{Cifar100-C}.}
	\label{fig:dataser_severity_levels}
\end{figure}

شکل \ref{fig:dataset_corruptions}انواع تخریب‌های اعمالی با بیشترین شدت در مجموعه داده‌های ذکر شده را نشان می‌دهد. در حالت انطباق برخط زمان آزمون هر کدام از تخریب‌ها به صورت برخط و جریان داده‌ای به مدل برای انطباق داده می شود و عملکرد مدل روی همان تخریب اندازه گیری می‌شود. در حالت انطباق برخط مستمر زمان آزمون هر کدام از تخریب‌ها به صورت جریان داده‌ای و پیوسته به مدل داده می‌شود و مدل بلافاصله بعد از انطباق بر روی یک تخریب با تخریب دیگری بدون دانستن زمان وقوع این تغییر مواجه می‌شود و می‌بایست فرایند انطباق را ادامه دهد. شکل \ref{fig:dataset_continual} نمایی از ترتیب مواجهه‌ی مدل با داده‌های مخدوش در حالت انطباق مستمر را نشان می‌دهد. همان‌طور که مشاهده می‌شود، مدل در طول جریان آزمون با یک توالی از نمونه‌های مخدوش مواجه می‌شود که تمامی آن‌ها با بیشترین شدت تخریب (سطح ۵) ارائه شده‌اند. این ترتیب به منظور بررسی پایداری و توانایی مدل در وفق یافتن با تغییرات توزیع داده‌ها طراحی شده و نقش کلیدی در ارزیابی عملکرد روش پیشنهادی در شرایط برخط و مستمر دارد. در واقع، در حالت انطباق برخط عادی، مدل برای هر تخریب به صورت مستقل پیش‌بینی و انطباق انجام می‌دهد و سپس قبل از مواجهه با تخریب بعدی، وزن‌های مدل به حالت اولیه و از پیش آموزش‌دیده بازنشانی می‌شوند. در مقابل، در حالت مستمر، مدل بدون اطلاع از زمان تغییر تخریب، وزن‌های خود را به صورت پیوسته به‌روزرسانی می‌کند و باید توانایی حفظ دقت و جلوگیری از انباشت خطا را در طول مواجهه با توالی تخریب‌ها نشان دهد.
علاوه بر این، طراحی توالی تخریب‌ها با بیشترین شدت (سطح ۵) به گونه‌ای انجام شده است که شرایطی چالش‌برانگیز برای مدل ایجاد شود و توانایی آن در مقابله با تغییرات شدید توزیع داده‌ها و حفظ عملکرد بلندمدت مورد سنجش قرار گیرد. این طراحی امکان بررسی دقیق جنبه‌های مختلف رفتار مدل از جمله پایداری پیش‌بینی‌ها در مواجهه با داده‌های نامطمئن، انعطاف‌پذیری مدل در تنظیم خود به تغییرات ناگهانی، و اثربخشی سازوکار‌های جلوگیری از انباشت خطا در طول زمان را فراهم می‌کند.
همچنین، این حالتها امکان مقایسه‌ی عملکرد مدل را در دو حالت متفاوت ارائه می‌دهند: حالت مستقل، که در آن مدل بعد از هر تخریب به وضعیت اولیه بازمی‌گردد، و حالت مستمر، که در آن مدل به‌صورت پیوسته وزن‌های خود را به‌روزرسانی می‌کند. این مقایسه، دیدی روشن از توانایی مدل در حفظ دقت و ثبات عملکرد در محیط‌های پویا و چالش‌برانگیز ارائه می‌دهد و کمک می‌کند نقاط ضعف احتمالی در مواجهه با داده‌های مخدوش و با شدت‌های مختلف تخریب شناسایی شود.


\begin{figure}[H]
	\centering
	\subfigure[{Glass}]{\includegraphics[width=0.18\linewidth]{images/severity/cifar100c/severity_5/glass_blur_17.png}}
	\subfigure[{\scriptsize Defoc.}]{\includegraphics[width=0.18\linewidth]{images/severity/cifar100c/severity_5/defocus_blur_17.png}}
	\subfigure[{Impul.}]{\includegraphics[width=0.18\linewidth]{images/severity/cifar100c/severity_5/impulse_noise_17.png}}
	\subfigure[{Shot}]{\includegraphics[width=0.18\linewidth]{images/severity/cifar100c/severity_5/shot_noise_17.png}}
	\subfigure[{Gauss.}]{\includegraphics[width=0.18\linewidth]{images/severity/cifar100c/severity_5/gaussian_noise_17.png}}
	
	\subfigure[{Fog}]{\includegraphics[width=0.18\linewidth]{images/severity/cifar100c/severity_5/fog_17.png}}
	\subfigure[{Frost}]{\includegraphics[width=0.18\linewidth]{images/severity/cifar100c/severity_5/frost_17.png}}
	\subfigure[{Snow}]{\includegraphics[width=0.18\linewidth]{images/severity/cifar100c/severity_5/snow_17.png}}
	\subfigure[{Zoom}]{\includegraphics[width=0.18\linewidth]{images/severity/cifar100c/severity_5/zoom_blur_17.png}}
	\subfigure[{Motion}]{\includegraphics[width=0.18\linewidth]{images/severity/cifar100c/severity_5/motion_blur_17.png}}
	
	\subfigure[{JPEG}]{\includegraphics[width=0.18\linewidth]{images/severity/cifar100c/severity_5/jpeg_compression_17.png}}
	\subfigure[{Pixel}]{\includegraphics[width=0.18\linewidth]{images/severity/cifar100c/severity_5/pixelate_17.png}}
	\subfigure[{Elastic}]{\includegraphics[width=0.18\linewidth]{images/severity/cifar100c/severity_5/elastic_transform_17.png}}
	\subfigure[{Contr.}]{\includegraphics[width=0.18\linewidth]{images/severity/cifar100c/severity_5/contrast_17.png}}
	\subfigure[{Brit.}]{\includegraphics[width=0.18\linewidth]{images/severity/cifar100c/severity_5/brightness_17.png}}
	
	\caption[انواع تخریب‌ داده‌های مورد بررسی در مجموعه داده‌های \lr{Cifar10-C,Cifar100-C}.]{انواع تخریب‌ داده‌های مورد بررسی در مجموعه داده‌های \lr{Cifar10-C,Cifar100-C}.}
	\label{fig:dataset_corruptions}
\end{figure}

\begin{figure}[H]
	\centering
	\subfigure[{\tiny \rotatebox{45}{JPEG}}]{\includegraphics[width=0.06\linewidth]{images/severity/cifar100c/severity_5/jpeg_compression_17.png}}
	\subfigure[{\tiny \rotatebox{45}{Pixel}}]{\includegraphics[width=0.06\linewidth]{images/severity/cifar100c/severity_5/pixelate_17.png}}
	\subfigure[{\tiny \rotatebox{45}{Elastic}}]{\includegraphics[width=0.06\linewidth]{images/severity/cifar100c/severity_5/elastic_transform_17.png}}
	\subfigure[{\tiny \rotatebox{45}{Contr.}}]{\includegraphics[width=0.06\linewidth]{images/severity/cifar100c/severity_5/contrast_17.png}}
	\subfigure[{\tiny \rotatebox{45}{Brit.}}]{\includegraphics[width=0.06\linewidth]{images/severity/cifar100c/severity_5/brightness_17.png}}
	\subfigure[{\tiny \rotatebox{45}{Fog}}]{\includegraphics[width=0.06\linewidth]{images/severity/cifar100c/severity_5/fog_17.png}}
	\subfigure[{\tiny \rotatebox{45}{Frost}}]{\includegraphics[width=0.06\linewidth]{images/severity/cifar100c/severity_5/frost_17.png}}
	\subfigure[{\tiny \rotatebox{45}{Snow}}]{\includegraphics[width=0.06\linewidth]{images/severity/cifar100c/severity_5/snow_17.png}}
	\subfigure[{\tiny \rotatebox{45}{Zoom}}]{\includegraphics[width=0.06\linewidth]{images/severity/cifar100c/severity_5/zoom_blur_17.png}}
	\subfigure[{\tiny \rotatebox{45}{Motion}}]{\includegraphics[width=0.06\linewidth]{images/severity/cifar100c/severity_5/motion_blur_17.png}}
	\subfigure[{\tiny \rotatebox{45}{Glass}}]{\includegraphics[width=0.06\linewidth]{images/severity/cifar100c/severity_5/glass_blur_17.png}}
	\subfigure[{\tiny \rotatebox{45}{Defoc.}}]{\includegraphics[width=0.06\linewidth]{images/severity/cifar100c/severity_5/defocus_blur_17.png}}
	\subfigure[{\tiny \rotatebox{45}{Impul.}}]{\includegraphics[width=0.06\linewidth]{images/severity/cifar100c/severity_5/impulse_noise_17.png}}
	\subfigure[{\tiny \rotatebox{45}{Shot}}]{\includegraphics[width=0.06\linewidth]{images/severity/cifar100c/severity_5/shot_noise_17.png}}
	\subfigure[{\tiny \rotatebox{45}{Gauss.}}]{\includegraphics[width=0.06\linewidth]{images/severity/cifar100c/severity_5/gaussian_noise_17.png}}
	\\[4pt]
	$\xrightarrow{\hspace{1\linewidth}}$
	\caption[ترتیب تخریب‌ها با بیشترین میزان شدت در انطباق در زمان آزمون مستمر]{ترتیب تخریب‌ها با بیشترین میزان شدت در انطباق در زمان آزمون مستمر.}
	\label{fig:dataset_continual}
\end{figure}
\newpage
\قسمت{معیارهای ارزیابی}

برای ارزیابی عملکرد روش پیشنهادی، معیار‌های زیر در نظر گرفته شده‌است که به شرح زیر توضیح داده می‌شوند:

\begin{itemize}
	\item \textbf{دقت (Accuracy):} 
	دقت، نسبت نمونه‌هایی است که مدل به درستی پیش‌بینی کرده است. در این پژوهش، دقت بر روی نمونه‌های مخدوش اندازه‌گیری می‌شود تا توانایی مدل در حفظ صحت پیش‌بینی‌ها در شرایط تغییر توزیع داده‌ها ارزیابی گردد.
	
	\item \textbf{خطای کالیبراسیون مورد انتظار \lr{(ECE)}:} 
	معیار \lr{ECE}\پانویس{Expected Calibration Error} میزان همخوانی احتمال پیش‌بینی شده توسط مدل با رخداد واقعی آن را سنجش می‌کند. به طور مثال، اگر مدل برای مجموعه‌ای از نمونه‌ها با اطمینان ۸۰\% پیش‌بینی ارائه دهد، در حالت ایده‌آل حدود ۸۰\% این نمونه‌ها باید به طور صحیح طبقه‌بندی شده باشند. 
	معیار ECE با تقسیم محدوده احتمال به $M$ بازه و محاسبه اختلاف میانگین پیش‌بینی و نسبت واقعی نمونه‌های درست در هر بازه به دست می‌آید و میانگین وزنی این اختلافات نشان‌دهنده ECE است. این معیار طبق رابطه‌ی  \ref{eq:ece} محاسبه می‌شود.
	\begin{equation}
		\label{eq:ece}
		\text{ECE} = \sum_{m=1}^{M} \frac{|B_m|}{n} \big| \text{acc}(B_m) - \text{conf}(B_m) \big|
	\end{equation}
	که $B_m$ مجموعه نمونه‌های با احتمال پیش‌بینی در بازه $m$، $|B_m|$ تعداد نمونه‌ها در این بازه، $n$ تعداد کل نمونه‌ها، $\text{acc}(B_m)$ دقت واقعی و $\text{conf}(B_m)$ میانگین احتمال پیش‌بینی در بازه $m$ است.
	
	\item \textbf{بیشینه خطای کالیبراسیون \lr{(MCE)}:} 
	معیار \lr{MCE}\پانویس{Maximum Calibration Error} بیشترین اختلاف بین احتمال پیش‌بینی شده و نسبت واقعی نمونه‌های درست در بین تمام بازه‌ها را اندازه‌گیری می‌کند. این معیار طبق رابطه‌ی \ref{eq:mce} محاسبه می شود.
	\begin{equation}
		\label{eq:mce}
		\text{MCE} = \max_{m \in \{1,\dots,M\}} \big| \text{acc}(B_m) - \text{conf}(B_m) \big|
	\end{equation}
	برخلاف ECE که میانگین خطا را نشان می‌دهد، MCE حساس به بدترین حالت کالیبراسیون مدل است و نقاطی که مدل بیشترین اشتباه پیش‌بینی را دارد مشخص می‌کند.
	
	\item \textbf{مساحت زیر منحنی مشخصه عملکرد گیرنده \lr{(AUROC)}:} 
	معیار \lr{AUROC}\پانویس{Area Under the Receiver Operating Characteristic} توانایی مدل در تمایز نمونه‌های درون توزیع از نمونه‌های خارج از توزیع را اندازه‌گیری می‌کند. در این پژوهش حین فرایند انطباق نمونه‌هایی که پیش‌بینی مدل با مقدار واقعی آن‌ها مطابقت دارد، به عنوان نمونه‌های درون توزیع و نمونه‌هایی که پیش‌بینی مدل با مقدار واقعی مطابقت ندارد، به عنوان نمونه‌های خارج از توزیع در نظر گرفته می‌شوند.
	
	از نظر الگوریتمی، AUROC تمام نمونه‌های درون توزیع را با تمام نمونه‌های خارج از توزیع جفت می‌کند. هر بار که امتیاز شناسایی خارج از توزیع داده‌ی درون توزیع، که در این پژوهش منفی امتیاز جرم تقریبی می‌باشد، از مقدار نمونه خارج از توزیع بیشتر باشد، یک واحد به شمارنده اضافه می‌شود و اگر مقادیر برابر باشند نیم واحد در نظر گرفته می‌شود. مجموع این مقادیر بر تعداد کل جفت‌ها تقسیم می‌شود تا شاخص AUROC به دست آید. این شاخص عددی بین صفر و یک است و هر چه به یک نزدیک‌تر باشد، نشان‌دهنده توانایی بالای مدل در تمایز نمونه‌های درست از نمونه‌های اشتباه است. این معیار طبق رابطه‌ی \ref{eq:auroc} محاسبه می‌شود.
	\begin{equation}
		\label{eq:auroc}
		\text{AUROC} = \frac{\sum_{i \in \text{ID}} \sum_{j \in \text{OOD}} \mathbb{1}\big( s_i > s_j \big) + \frac{1}{2}\mathbb{1}\big( s_i = s_j \big)}{|\text{ID}| \cdot |\text{OOD}|}
	\end{equation}
	که در آن $s_i$ و $s_j$ مقادیر منفی امتیاز جرم تقریبی برای نمونه‌های درون و خارج از توزیع هستند و $\mathbb{1}$ تابع شاخص است.
\end{itemize}

\قسمت{نحوه‌ی پیاده‌سازی و ارزیابی}

برای پیاده‌سازی روش پیشنهادی، از مدل‌های مبتنی بر مبدل بینایی استفاده شده است. در این پیاده‌سازی، انطباق تنها بر روی لایه های نرمال‌سازی لایه‌ای اعمال شده است و همانطور که پیش‌تر نیز بیان شد نسبت به اندازه‌ی دسته حساسیت نداشته و به‌طور کلی با پارامتر‌های کم برای مسئله‌ی انطباق زمان آزمون برخط گزینه‌ای ایده‌آل می‌باشد، پارامتر‌های این لایه ها سهم بسیار اندکی از کل وزن‌های شبکه، کمتر از $0.05\%$، را تشکیل می‌دهد. این طراحی، ضمن حفظ ساختار اصلی مدل و جلوگیری از افزایش بیش‌ازحد پارامترهای قابل آموزش، امکان به‌روزرسانی سریع و پایدار در شرایط برخط را فراهم می‌کند.  

تمامی آزمایش‌ها با استفاده از مدل‌های از پیش آموزش دیده در \lr{HuggingFace} انجام شده‌است تا بتوان از وزن‌های پایه‌ای قوی بهره برد و نتایج واقعی عملکرد روش پیشنهادی را در شرایط انطباق برخط و مستمر ارزیابی نمود. استفاده از مدل‌های پیش‌آموزش دیده همچنین به کاهش زمان آموزش و بهبود همگرایی کمک می‌کند. این موارد در قسمت پیوست آورده شده‌است.

مسئله در دو حالت اصلی انطباق برخط زمان آزمون و انطباق برخط مستمر زمان آزمون ارزیابی شده است.در هر دو حالت، تنظیمات مشابهی از جمله اندازه‌ی دسته برای تمامی ارزیابی‌ها برابر با ۳۲ در نظر گرفته شده است تا مقایسه‌ی عملکرد روش‌ها به صورت منسجم و یکنواخت انجام شود.  

برای مقایسه‌ی روش پیشنهادی، روش‌های پایه‌ای \lr{TENT}، \lr{SAR} و \lr{TEA} در نظر گرفته شده‌اند. علاوه بر این، معیار امتیاز جرم تقریبی به هر یک از روش‌های پایه‌ای تزریق شده و نتایج در دو حالت استاندارد و با اعمال امتیاز جرم تقریبی، گزارش شده است.  

لازم به ذکر است که پیش از اعمال هر نوع معیاری در فرآیند انطباق زمان آزمون، لازم است اعتبار آن معیار از منظر ارتباط با عملکرد مدل بررسی شود. به بیان دیگر، باید تحلیل شود که آیا بین مقادیر معیار مورد استفاده مانند آنتروپی، انرژی یا امتیاز جرم تقریبی و نرخ خطای مدل، همبستگی معناداری وجود دارد یا خیر. تنها در صورتی که چنین همبستگی‌ای برقرار باشد، معیار می‌تواند نقش مؤثری در هدایت فرآیند انطباق ایفا کند. به همین دلیل، در ابتدای این فصل یک مرحله تحلیل اعتبار برای معیارهای مورد نظر بر روی مجموعه داده‌های تحت بررسی انجام شده است تا پایه‌ی محکمی برای ارزیابی‌های بعدی فراهم گردد.

به منظور انتخاب بهترین مجموعه فراپارامترها، الگوریتم \lr{Grid Search} به کار گرفته شده است. در این فرآیند، هم نتایج حاصل از بهترین ترکیب پارامترها و هم میانگین عملکرد تمام ترکیب‌های مورد آزمون گزارش شده‌اند. توجه به میانگین نتایج و واریانس آن اهمیت ویژه‌ای دارد، زیرا در حالت انطباق برخط زمان آزمون عملاً امکان تنظیم دقیق پارامترها وجود ندارد و مدل باید بتواند با مجموعه‌ای از پارامترهای تقریباً بهینه، عملکرد پایدار و قابل اعتمادی ارائه دهد.

برای تحلیل حساسیت و استحکام مدل نسبت به تغییرات فراپارامترها، نمودارهای سنجش حساسیت و استحکام ارائه شده‌اند که نشان می‌دهند مدل چگونه در مواجهه با مقادیر مختلف پارامترها دقت و پایداری خود را حفظ می‌کند. این نمودارها کمک می‌کنند تا انعطاف‌پذیری، استحکام و حساسیت روش در شرایط واقعی که عملا امکان تنظیم دقیق فراپارامتر‌های شبکه وجود ندارد ارزیابی شود.

برای تحلیل بصری تأثیر تغییرات بر روی تصمیمات مدل، بصری‌سازی \lr{Grad-CAM}\cite{gradcam} نیز ارائه شده است. این ابزار امکان مشاهده‌ی نواحی فعال در تصویر که بیشترین اثر را بر پیش‌بینی مدل دارند فراهم می‌کند و به فهم بهتر فرآیند انطباق و چگونگی پاسخ مدل به داده‌های مخدوش کمک می‌کند. به این ترتیب، ترکیب تحلیل کمی و کیفی، تصویری جامع از عملکرد روش پیشنهادی در حالتهای برخط و مستمر ارائه می‌دهد.

\زیرقسمت{تحلیل اعتبار}

پیش از آنکه معیارهای مختلف برای هدایت فرآیند انطباق زمان آزمون به کار گرفته شوند، لازم است اعتبار آن‌ها از منظر ارتباط با عملکرد مدل مورد بررسی قرار گیرد. به بیان دقیق‌تر، معیار مورد استفاده باید نشان دهد که تغییرات آن به‌طور معناداری با نرخ خطای مدل همبسته است. وجود چنین همبستگی‌ای تضمین می‌کند که معیار می‌تواند به‌عنوان یک سیگنال معتبر برای تصمیم‌گیری در فرآیند انطباق به کار گرفته شود؛ در غیر این صورت، استفاده از آن می‌تواند باعث بروز انطباق نادرست و نتایج ناخواسته شود.

برای بررسی این موضوع، یک تحلیل همبستگی میان مقادیر معیار و نرخ خطای مدل در شرایط تخریب‌های مختلف انجام شده است. نتایج این تحلیل در شکل \ref{fig:validity_analysis} برای مجموعه داده‌های \lr{CIFAR10-C} و \lr{CIFAR100-C} ارائه شده است. نتایج به‌دست‌آمده نشان می‌دهند که میان مقادیر معیارهای بررسی‌شده و نرخ خطای مدل یک همبستگی مستقیم وجود دارد؛ به این معنا که با افزایش خطا، مقدار معیار نیز افزایش می‌یابد. این همبستگی بیانگر آن است که معیار انتخابی می‌تواند رفتار خطای مدل را به‌طور معناداری بازتاب دهد و در نتیجه به‌عنوان یک سیگنال معتبر برای هدایت فرآیند انطباق زمان آزمون به کار گرفته شود.

به طور کلی، این تحلیل اعتبار به‌عنوان گام مقدماتی، مبنای محکمی برای استفاده از معیارها در ادامه‌ی آزمایش‌ها فراهم می‌سازد و تضمین می‌کند که نتایج حاصل از اعمال آن‌ها بر فرآیند انطباق، دارای پشتوانه‌ی مناسب باشند.

\begin{figure}[t]
	\centering
	% ---- CIFAR10-C ----
	\subfigure[تحلیل اعتبار در مجموعه داده‌ی \lr{CIFAR10-C}.]{
		\begin{minipage}{0.9\linewidth}
			\centering
			\subfigure[آنتروپی]{
				\includegraphics[width=0.3\linewidth]{images/error_entropy_all(cifar10c).png}
			}
			\subfigure[انرژی]{
				\includegraphics[width=0.3\linewidth]{images/error_energy_all(cifar10c).png}
			}
			\subfigure[امتیاز جرم تقریبی]{
				\includegraphics[width=0.3\linewidth]{images/error_approx_mass_all(cifar10c).png}
			}
		\end{minipage}
	}
	% ---- CIFAR100-C ----
	\subfigure[تحلیل اعتبار در مجموعه داده‌ی \lr{CIFAR100-C}.]{
		\begin{minipage}{0.9\linewidth}
			\centering
			\subfigure[آنتروپی]{
				\includegraphics[width=0.3\linewidth]{images/error_entropy_all(cifar100c).png}
			}
			\subfigure[انرژي]{
				\includegraphics[width=0.3\linewidth]{images/error_energy_all(cifar100c).png}
			}
			\subfigure[امتیاز جرم تقریبی]{
				\includegraphics[width=0.3\linewidth]{images/error_approx_mass_all(cifar100c).png}
			}
		\end{minipage}
	}
	\caption[تحلیل اعتبار معیار]{تحلیل اعتبار معیارهای آنتروپی، انرژی و امتیاز جرم تقریبی با بررسی همبستگی آن‌ها با نرخ خطای مدل در مجموعه داده‌های \lr{CIFAR10-C} و \lr{CIFAR100-C}. محور افقی مقادیر معیار نرمال‌شده در بازه‌های مختلف و محور عمودی نرخ خطای مدل می‌باشند.}
	
	\label{fig:validity_analysis}
\end{figure}

\زیرقسمت{نتایج}

در این بخش، عملکرد روش پیشنهادی در دو حالت اصلی انطباق برخط زمان آزمون و انطباق برخط مستمر زمان آزمون مورد ارزیابی قرار گرفته است. در هر دو حالت، تنظیمات آزمایش‌ها یکسان نگه داشته شده‌اند تا امکان مقایسه‌ای منسجم و دقیق میان روش‌های مختلف فراهم شود. این هماهنگی در تنظیمات، موجب می‌شود هر تفاوتی در نتایج مستقیماً ناشی از کارایی روش‌ها باشد و اثر عوامل جانبی کاهش یابد.

روش‌های پایه‌ای \lr{TENT}، \lr{SAR} و \lr{TEA} به عنوان معیارهای مرجع در نظر گرفته شده‌اند و در هر یک از این روش‌ها، اثر تزریق معیار امتیاز جرم تقریبی نیز بررسی شده است. ارائه‌ی نتایج در دو حالت استاندارد و با اعمال معیار پیشنهادی امکان می‌دهد تا نقش و تأثیر این معیار در هدایت فرآیند انطباق زمان آزمون به‌وضوح مشخص شود و نشان دهد چگونه می‌تواند دقت و پایداری مدل را ارتقا دهد.

هدف اصلی از ارائه‌ی این نتایج، نه تنها نمایش میزان بهبود عملکرد با بهره‌گیری از معیار مذکور است، بلکه سنجش توانایی روش پیشنهادی در حفظ ثبات و دقت در شرایط متغیر و چالش‌برانگیز داده‌ها نیز مدنظر قرار گرفته است.

با توجه به اینکه در مسئله‌ی انطباق برخط زمان آزمون امکان تنظیم دقیق فراپارامترها وجود ندارد، نتایج میانگین عملکرد تمام ترکیب‌های پارامترها که در فرآیند \lr{Grid Search} بررسی شده‌اند نیز ارائه می‌شوند. این رویکرد نشان می‌دهد که روش پیشنهادی علاوه بر دستیابی به بهترین عملکرد، در شرایطی که پارامترها به صورت تقریباً بهینه انتخاب شده‌اند نیز پایداری و قابلیت اطمینان خود را حفظ می‌کند و واریانس پایینی دارد.

در ادامه، جداول \ref{table:t1} تا \ref{table:t8} ارائه شده‌اند که نتایج حاصل از ارزیابی روش پیشنهادی و روش‌های پایه‌ای را در حالتهای مختلف انطباق برخط آزمون و انطباق برخط آزمون مستمر نشان می‌دهند. این جداول امکان مقایسه‌ی کمی دقیق میان روش‌ها را فراهم می‌کنند و اثر تزریق معیار امتیاز جرم تقریبی و حساسیت روش‌ها به انتخاب پارامترها را نیز به نمایش می‌گذارند.
همچنین جدول \ref{table:alg2_results} نتایج اجرای الگوریتم \ref{alg:layer_selection} را نشان می‌دهد. لایه‌ی انتخاب‌شده توسط این الگوریتم به الگوریتم \ref{alg:ocetta} انتقال یافته است تا امتیاز جرم تقریبی بر خروجی آن لایه اعمال گردد. مطابق جدول \ref{table:t9}، این رویکرد اگرچه نسبت به مدل بدون انطباق منبع بهبودهایی را نشان می‌دهد، اما عملکرد نسخه‌ی اصلی یعنی الگوریتم \ref{alg:ocetta} هم‌چنان اندکی بهتر است. این یافته نشان می‌دهد که مسیرهای پژوهشی آینده می‌توانند بر توسعه‌ی رویکردهای جایگزین برای تولید داده‌های مصنوعی خارج از توزیع و به‌کارگیری روش‌های متنوع در حضور داده‌های منبع متمرکز شوند، تا بدین ترتیب عملکرد روش پیشنهادی نسبت به نسخه‌ی اصلی الگوریتم بهبود یابد.\\

\begin{landscape}
\begin{table}[H]
    \centering
    \large
    \caption[مقایسه‌ی عملکرد در مجموعه داده‌ی \lr{CIFAR10-C}]{مقایسه‌ی عملکرد بر روی مدل ViT در مجموعه داده‌ی \lr{CIFAR10-C}}
    \label{table:t1}
    \begin{adjustbox}{width=1.5\textwidth}
    \begin{latin}
    \begin{tabular}{lccccccccccccccc|cccc}
    \toprule
    \multirow{2}{*}{Method} & \multicolumn{3}{c}{Noise} & \multicolumn{4}{c}{Blur} & \multicolumn{3}{c}{Weather} & \multicolumn{5}{c}{Digital} & \multicolumn{4}{c}{Avg} \\
    \cmidrule(lr){2-4} \cmidrule(lr){5-8}  \cmidrule(lr){9-11} \cmidrule(lr){12-16} \cmidrule(lr){17-20}
    & Gauss. & Shot & Impul. & Defoc. & Glass & Motion & Zoom & Snow & Frost & Fog & Brit. & Contr. & Elastic & Pixel & JPEG & Acc (↑) & AUROC (↑) & ECE (↓) &  MCE (↓)  \\
    \midrule
    Source & 76.23 & 77.04 & 74.64 & 92.84 & 69.25 & 90.18 & 94.00 & 93.72 & 92.18 & 90.21 & 97.69 & 92.99 & 82.88 & 57.22 & 84.45 & 84.37 & 87.84 & 9.52 & 34.85 \\
    TENT\cite{ftent} & 84.10 & 80.59 & 82.86 & \textbf{94.61} & 81.39 & 88.90 & 95.45 & 94.03 & 94.69 & 92.19 & 97.40 & 95.36 & 85.41 & 86.42 & 84.99 & 89.23 & 87.87 & 8.39 & 44.83 \\
    SAR\cite{sar} & 82.43 & 82.61 & 82.49 & 93.89 & 77.87 & 92.52 & 95.05 & 94.20 & 94.29 & 93.05 & \textbf{97.87} & 95.20 & 88.07 & 86.11 & 83.75 & 89.29 & 88.58 & 7.35 & 33.64 \\
    TEA\cite{tea} & 78.33 & 84.36 & 83.65 & 94.52 & 81.54 & 93.20 & \textbf{95.83} & 94.71 & 94.91 & 93.42 & 97.79 & 95.83 & 88.55 & 86.84 & 85.31 & 89.92 & 88.07 & 7.37 & 32.94 \\
    TENT(APMS) & \textbf{85.54} & 87.02 & \textbf{86.06} & 94.42 & \textbf{84.62} & \textbf{93.54} & 95.72 & 94.85 & 94.92 & \textbf{93.77} & 97.39 & 96.13 & \textbf{89.64} & \textbf{93.06} & \textbf{87.97} & \textbf{91.64} & 88.65 & \textbf{6.75} & 36.17 \\
    SAR(APMS) & 84.48 & \textbf{87.21} & 85.99 & 94.49 & 81.74 & 93.30 & 95.74 & \textbf{94.86} & \textbf{95.11} & 93.35 & 97.78 & \textbf{96.21} & 89.11 & 91.46 & 86.90 & 91.18 & \textbf{89.67} & 6.87 & 36.83 \\
    TEA(APMS) & 81.94 & 83.83 & 84.12 & 94.03 & 77.88 & 92.56 & 95.31 & 94.22 & 94.30 & 92.69 & 97.55 & 95.47 & 87.55 & 86.64 & 83.60 & 89.45 & 88.35 & 7.65 & \textbf{32.11} \\
    \end{tabular}
    \end{latin}
    \end{adjustbox}
\end{table}

\begin{table}[H]
    \centering
    \caption[مقایسه‌ی میانگین عملکرد در مجموعه داده‌ی \lr{CIFAR10-C}]{مقایسه‌ی میانگین عملکرد بر روی مدل ViT در مجموعه داده‌ی \lr{CIFAR10-C}}
    \label{table:t2}
    \begin{adjustbox}{width=1.5\textwidth}
    \begin{latin}
    \begin{tabular}{lccccccccccccccc|cccc}
    \toprule
    \multirow{2}{*}{Method} & \multicolumn{3}{c}{Noise} & \multicolumn{4}{c}{Blur} & \multicolumn{3}{c}{Weather} & \multicolumn{5}{c}{Digital} & \multicolumn{4}{c}{Avg} \\
    \cmidrule(lr){2-4} \cmidrule(lr){5-8}  \cmidrule(lr){9-11} \cmidrule(lr){12-16} \cmidrule(lr){17-20}
    & Gauss. & Shot & Impul. & Defoc. & Glass & Motion & Zoom & Snow & Frost & Fog & Brit. & Contr. & Elastic & Pixel & JPEG & Acc (↑) & AUROC (↑) & ECE (↓) &  MCE (↓)  \\
    \midrule
    Source & 76.23 & 77.04 & 74.64 & 92.84 & 69.25 & 90.18 & 94.00 & 93.72 & 92.18 & 90.21 & 97.69 & 92.99 & 82.88 & 57.22 & 84.45 & 84.37 & 80.12 & 9.52 & \textbf{34.85} \\
    TENT\cite{ftent} & 71.23 & 64.10 & 64.20 & 91.87 & 69.51 & 90.01 & 85.74 & 84.61 & 89.52 & 85.28 & 96.89 & 94.04 & 77.32 & 67.91 & 62.89 & 79.67 & 82.92 & 17.82 & 44.31 \\
    SAR\cite{sar} & 61.19 & 62.81 & 61.99 & 73.69 & 57.12 & 72.32 & 74.74 & 74.01 & 73.77 & 72.13 & 77.74 & 74.41 & 67.94 & 62.90 & 66.59 & 68.89 & 68.45 & 25.64 & 46.00 \\
    TEA\cite{tea} & 48.98 & 55.21 & 64.46 & 82.38 & 55.14 & 87.44 & 85.49 & 94.21 & 94.01 & 77.37 & 97.47 & 87.10 & 72.46 & 59.37 & 59.64 & 74.72 & 80.85 & 23.28 & 49.96 \\
    TENT(APMS) & \textbf{83.93} & \textbf{85.93} & \textbf{84.94} & \textbf{94.42} & \textbf{80.62} & \textbf{93.25} & \textbf{95.70} & \textbf{94.77} & \textbf{94.85} & \textbf{93.57} & 97.54 & \textbf{96.00} & \textbf{89.12} & \textbf{90.62} & \textbf{86.76} & \textbf{90.80} & \textbf{90.26} & \textbf{7.08} & 38.53 \\
    SAR(APMS) & 76.75 & 81.37 & 79.04 & 93.40 & 74.39 & 91.35 & 94.59 & 94.02 & 93.17 & 91.27 & \textbf{97.70} & 94.27 & 85.00 & 72.60 & 84.52 & 86.90 & 74.64 & 8.67 & 38.03 \\
    TEA(APMS) & 54.71 & 57.21 & 62.34 & 81.52 & 49.45 & 72.12 & 89.78 & 84.75 & 85.15 & 77.01 & 96.99 & 92.48 & 64.02 & 75.19 & 57.55 & 73.35 & 82.05 & 24.54 & 52.09 \\
    \end{tabular}
    \end{latin}
    \end{adjustbox}
\end{table}
\end{landscape}

\begin{landscape}
\begin{table}[H]
    \centering
    \caption[مقایسه‌ی عملکرد در مجموعه داده‌ی \lr{CIFAR10-C}]{مقایسه‌ی عملکرد بر روی مدل ViT در مجموعه داده‌ی \lr{CIFAR100-C}}
    \label{table:t3}
    \begin{adjustbox}{width=1.5\textwidth}
    \begin{latin}
    \begin{tabular}{lccccccccccccccc|cccc}
    \toprule
    \multirow{2}{*}{Method} & \multicolumn{3}{c}{Noise} & \multicolumn{4}{c}{Blur} & \multicolumn{3}{c}{Weather} & \multicolumn{5}{c}{Digital} & \multicolumn{4}{c}{Avg} \\
    \cmidrule(lr){2-4} \cmidrule(lr){5-8}  \cmidrule(lr){9-11} \cmidrule(lr){12-16} \cmidrule(lr){17-20}
    & Gauss. & Shot & Impul. & Defoc. & Glass & Motion & Zoom & Snow & Frost & Fog & Brit. & Contr. & Elastic & Pixel & JPEG & Acc (↑) & AUROC (↑) & ECE (↓) &  MCE (↓)  \\
    \midrule
    Source & 44.99 & 46.80 & 47.25 & 82.05 & 44.04 & 76.54 & 84.05 & 78.80 & 75.59 & 67.81 & 87.67 & 63.71 & 65.82 & 55.31 & 64.98 & 65.69 & 80.79 & 16.16 & 37.95 \\
    TENT\cite{ftent} & 48.26 & 40.07 & 54.53 & 84.12 & 53.47 & 80.23 & 85.13 & 80.41 & 78.88 & 71.79 & 88.33 & 75.66 & 72.25 & 76.03 & 66.71 & 70.39 & 81.86 & 18.99 & 39.89 \\
    SAR\cite{sar} & 59.50 & 62.55 & 62.11 & 83.86 & 62.74 & 79.29 & 84.91 & 80.10 & 78.66 & 73.35 & 88.01 & 76.92 & 72.47 & 77.36 & 67.22 & 73.94 & 38.37 & 13.93 & \textbf{34.93} \\
    TEA\cite{tea} & 59.47 & 37.25 & 65.65 & 84.55 & 65.10 & 81.51 & 85.91 & 81.35 & 80.59 & 75.65 & 88.77 & 80.23 & 74.50 & 79.15 & 68.53 & 73.88 & 76.18 & 16.35 & 37.42 \\
    TENT(APMS) & 67.04 & \textbf{69.70} & \textbf{72.32} & \textbf{85.04} & \textbf{71.21} & \textbf{83.42} & \textbf{86.68} & \textbf{83.24} & \textbf{83.57} & \textbf{79.45} & \textbf{89.24} & \textbf{85.22} & \textbf{77.68} & \textbf{83.90} & \textbf{73.29} & \textbf{79.40} & \textbf{87.04} & \textbf{13.50} & 37.12 \\
    SAR(APMS) & \textbf{67.45} & 69.43 & 71.32 & 84.58 & 69.80 & 82.15 & 86.23 & 81.78 & 82.08 & 76.98 & 88.92 & 81.65 & 76.20 & 82.60 & 71.72 & 78.19 & 48.05 & 13.60 & 39.87 \\
    TEA(APMS) & 66.57 & 69.27 & 71.98 & 83.99 & 70.18 & 82.85 & 86.18 & 82.37 & 82.65 & 78.82 & 88.55 & 84.31 & 76.94 & 83.35 & 72.24 & 78.68 & 86.35 & 13.84 & 39.96 \\
    \end{tabular}
    \end{latin}
    \end{adjustbox}
\end{table}

\begin{table}[H]
    \centering
    \caption[مقایسه‌ی میانگین عملکرد در مجموعه داده‌ی \lr{CIFAR100-C}]{مقایسه‌ی میانگین عملکرد بر روی مدل ViT در مجموعه داده‌ی \lr{CIFAR100-C}}
    \label{table:t4}
    \begin{adjustbox}{width=1.5\textwidth}
    \begin{latin}
    \begin{tabular}{lccccccccccccccc|cccc}
    \toprule
    \multirow{2}{*}{Method} & \multicolumn{3}{c}{Noise} & \multicolumn{4}{c}{Blur} & \multicolumn{3}{c}{Weather} & \multicolumn{5}{c}{Digital} & \multicolumn{4}{c}{Avg} \\
    \cmidrule(lr){2-4} \cmidrule(lr){5-8}  \cmidrule(lr){9-11} \cmidrule(lr){12-16} \cmidrule(lr){17-20}
    & Gauss. & Shot & Impul. & Defoc. & Glass & Motion & Zoom & Snow & Frost & Fog & Brit. & Contr. & Elastic & Pixel & JPEG & Acc (↑) & AUROC (↑) & ECE (↓) &  MCE (↓)  \\
    \midrule
    Source & 44.99 & 46.80 & 47.25 & 82.05 & 44.04 & 76.54 & 84.05 & 78.80 & 75.59 & 67.81 & 87.67 & 63.71 & 65.82 & 55.31 & 64.98 & 65.69 & 80.79 & 16.16 & 37.95 \\
    TENT\cite{ftent} & 36.54 & 25.13 & 28.89 & 76.45 & 32.43 & 74.71 & 84.16 & 70.85 & 71.47 & 58.97 & 88.13 & 77.92 & 57.47 & 67.08 & 48.18 & 59.89 & 75.72 & 32.36 & 50.02 \\
    SAR\cite{sar} & 37.27 & 28.06 & 42.42 & 64.00 & 45.30 & 60.42 & 65.58 & 61.12 & 61.13 & 55.38 & 68.43 & 57.98 & 54.30 & 58.57 & 50.98 & 54.06 & 50.87 & 22.14 & 45.02 \\
    TEA\cite{tea} & 40.48 & 27.80 & 43.44 & 78.86 & 58.67 & 81.38 & 85.66 & 81.10 & 79.57 & 75.92 & 88.52 & 81.30 & 71.83 & 80.18 & 61.27 & 69.07 & 75.80 & 23.22 & 45.79 \\
    TENT(APMS) & \textbf{63.63} & \textbf{66.40} & \textbf{68.78} & \textbf{84.71} & \textbf{67.07} & \textbf{82.39} & \textbf{86.35} & \textbf{82.17} & \textbf{82.43} & \textbf{77.72} & \textbf{89.01} & \textbf{82.21} & \textbf{75.77} & \textbf{81.58} & \textbf{71.21} & \textbf{77.43} & \textbf{86.05} & \textbf{13.68} & \textbf{36.89} \\
    SAR(APMS) & 50.70 & 46.85 & 47.06 & 72.54 & 44.77 & 68.10 & 74.06 & 69.60 & 67.62 & 61.44 & 77.10 & 59.53 & 59.62 & 59.42 & 57.91 & 61.09 & 48.23 & 14.10 & 40.90 \\
    TEA(APMS) & 55.63 & 59.62 & 61.16 & 83.73 & 60.08 & 81.26 & 85.60 & 81.14 & 81.14 & 76.53 & 88.29 & 80.55 & 68.93 & 79.71 & 62.76 & 73.74 & 84.65 & 16.85 & 44.23 \\
    \end{tabular}
    \end{latin}
    \end{adjustbox}
\end{table}
\end{landscape}

\begin{landscape}
\begin{table}[H]
    \centering
    \caption[مقایسه‌ی عملکرد در مجموعه داده‌ی \lr{CIFAR10-C} در حالت مستمر]{مقایسه‌ی عملکرد بر روی مدل ViT در مجموعه داده‌ی \lr{CIFAR10-C} در حالت مستمر}
    \label{table:t5}
    \begin{adjustbox}{width=1.5\textwidth}
    \begin{latin}
    \begin{tabular}{lccccccccccccccc|cccc}
    \toprule
    Time & \multicolumn{15}{l|}{$t \xrightarrow{\hspace*{23cm}}$} & \\ 
    \midrule
    \multirow{2}{*}{Method} & \multicolumn{3}{c}{Noise} & \multicolumn{4}{c}{Blur} & \multicolumn{3}{c}{Weather} & \multicolumn{5}{c}{Digital} & \multicolumn{4}{c}{Avg} \\
    \cmidrule(lr){2-4} \cmidrule(lr){5-8}  \cmidrule(lr){9-11} \cmidrule(lr){12-16} \cmidrule(lr){17-20}
    & Gauss. & Shot & Impul. & Defoc. & Glass & Motion & Zoom & Snow & Frost & Fog & Brit. & Contr. & Elastic & Pixel & JPEG & Acc (↑) & AUROC (↑) & ECE (↓) &  MCE (↓)  \\
    \midrule
    Source & 76.23 & 77.04 & 74.64 & 92.84 & 69.25 & 90.18 & 94.00 & 93.72 & 92.18 & 90.21 & \textbf{97.69} & 92.99 & 82.88 & 57.22 & 84.45 & 84.37 & 80.12 & 9.52 & 34.85 \\
    TENT\cite{ftent} & 79.61 & 83.19 & 68.21 & 89.98 & 74.59 & 87.37 & 91.66 & 91.40 & 93.80 & 89.98 & 96.73 & 92.21 & 86.71 & 78.71 & 85.23 & 85.96 & 85.07 & 11.94 & 39.49 \\
    SAR\cite{sar} & 81.51 & 86.19 & 80.87 & 93.14 & 80.38 & 90.36 & 93.88 & 93.57 & 94.67 & 91.58 & 96.76 & 93.12 & 87.45 & 86.18 & 85.85 & 89.03 & 71.08 & 8.35 & 35.17 \\
    TEA\cite{tea} & 78.51 & 82.40 & 79.33 & 93.47 & 77.55 & \textbf{92.29} & \textbf{95.45} & 94.04 & 94.36 & 92.33 & 96.90 & \textbf{94.46} & 87.18 & 54.28 & 80.72 & 86.22 & 87.92 & 10.70 & 43.84 \\
    TENT(APMS) & \textbf{83.17} & \textbf{87.48} & \textbf{85.75} & \textbf{93.83} & \textbf{82.45} & 91.69 & 94.89 & 93.75 & \textbf{95.03} & \textbf{92.45} & 96.98 & 94.02 & \textbf{88.18} & \textbf{90.02} & \textbf{86.82} & \textbf{90.43} & \textbf{89.85} & \textbf{8.29} & 45.52 \\
    SAR(APMS) & 79.30 & 82.50 & 78.48 & 93.30 & 77.57 & 91.67 & 94.92 & \textbf{94.26} & 94.47 & 92.13 & 96.91 & 94.40 & 87.09 & 79.77 & 86.05 & 88.19 & 67.16 & 8.57 & \textbf{33.32} \\
    TEA(APMS) & 78.55 & 82.34 & 80.54 & 93.66 & 75.18 & 91.99 & 95.09 & 93.48 & 93.68 & 92.02 & 96.10 & 93.34 & 84.96 & 51.25 & 77.91 & 85.34 & 87.66 & 11.42 & 40.93 \\
    \end{tabular}
    \end{latin}
    \end{adjustbox}
\end{table}

\begin{table}[H]
    \centering
    \caption[مقایسه‌ی میانگین عملکرد در مجموعه داده‌ی \lr{CIFAR10-C} در حالت مستمر]{مقایسه‌ی میانگین عملکرد بر روی مدل ViT در مجموعه داده‌ی \lr{CIFAR10-C} در حالت مستمر}
    \label{table:t6}
    \label{table:}
    \begin{adjustbox}{width=1.5\textwidth}
    \begin{latin}
    \begin{tabular}{lccccccccccccccc|cccc}
    \toprule
    Time & \multicolumn{15}{l|}{$t \xrightarrow{\hspace*{23cm}}$} & \\ 
    \midrule
    \multirow{2}{*}{Method} & \multicolumn{3}{c}{Noise} & \multicolumn{4}{c}{Blur} & \multicolumn{3}{c}{Weather} & \multicolumn{5}{c}{Digital} & \multicolumn{4}{c}{Avg} \\
    \cmidrule(lr){2-4} \cmidrule(lr){5-8}  \cmidrule(lr){9-11} \cmidrule(lr){12-16} \cmidrule(lr){17-20}
    & Gauss. & Shot & Impul. & Defoc. & Glass & Motion & Zoom & Snow & Frost & Fog & Brit. & Contr. & Elastic & Pixel & JPEG & Acc (↑) & AUROC (↑) & ECE (↓) &  MCE (↓)  \\
    \midrule
    Source & 76.23 & 77.04 & 74.64 & \textbf{92.84} & \textbf{69.25} & \textbf{90.18} & \textbf{94.00} & \textbf{93.72} & \textbf{92.18} & \textbf{90.21} & \textbf{97.69} & \textbf{92.99} & \textbf{82.88} & 57.22 & \textbf{84.45} & \textbf{84.37} & 80.12 & \textbf{9.52} & \textbf{34.85} \\
    TENT\cite{ftent} & 71.23 & 57.71 & 37.41 & 45.81 & 36.22 & 41.75 & 42.75 & 42.66 & 43.22 & 41.00 & 44.47 & 41.72 & 39.45 & 37.03 & 39.33 & 44.12 & 61.78 & 54.95 & 67.90 \\
    SAR\cite{sar} & 61.19 & 63.47 & 56.44 & 69.79 & 52.35 & 58.99 & 60.81 & 61.07 & 61.85 & 53.22 & 62.78 & 60.08 & 53.05 & 51.14 & 51.40 & 58.51 & 64.27 & 39.53 & 55.59 \\
    TEA\cite{tea} & 54.89 & 40.79 & 30.83 & 27.62 & 23.51 & 26.47 & 27.09 & 26.81 & 26.87 & 26.47 & 27.38 & 26.89 & 25.44 & 18.86 & 24.14 & 28.94 & 60.24 & 70.23 & 77.79 \\
    TENT(APMS) & \textbf{83.93} & \textbf{86.76} & \textbf{81.60} & 76.92 & 67.89 & 74.93 & 77.69 & 76.68 & 77.66 & 75.77 & 79.16 & 77.21 & 71.28 & \textbf{73.14} & 70.09 & 76.71 & \textbf{82.20} & 22.33 & 50.53 \\
    SAR(APMS) & 76.75 & 72.80 & 70.09 & 82.73 & 63.54 & 71.08 & 73.54 & 73.11 & 72.95 & 71.32 & 75.39 & 73.24 & 67.05 & 55.46 & 66.61 & 71.04 & 69.45 & 25.78 & 46.05 \\
    TEA(APMS) & 59.48 & 50.31 & 38.82 & 43.14 & 27.31 & 26.53 & 27.02 & 26.70 & 26.74 & 26.40 & 27.22 & 26.67 & 24.99 & 18.25 & 23.58 & 31.54 & 62.18 & 67.50 & 75.20 \\
    \end{tabular}
    \end{latin}
    \end{adjustbox}
\end{table}
\end{landscape}

\begin{landscape}
\begin{table}[H]
    \centering
    \caption[مقایسه‌ی عملکرد در مجموعه داده‌ی \lr{CIFAR100-C} در حالت مستمر]{مقایسه‌ی عملکرد بر روی مدل ViT در مجموعه داده‌ی \lr{CIFAR100-C} در حالت مستمر}
    \label{table:t7}
    \begin{adjustbox}{width=1.5\textwidth}
    \begin{latin}
    \begin{tabular}{lccccccccccccccc|cccc}
    \toprule
    Time & \multicolumn{15}{l|}{$t \xrightarrow{\hspace*{23cm}}$} & \\ 
    \midrule
    \multirow{2}{*}{Method} & \multicolumn{3}{c}{Noise} & \multicolumn{4}{c}{Blur} & \multicolumn{3}{c}{Weather} & \multicolumn{5}{c}{Digital} & \multicolumn{4}{c}{Avg} \\
    \cmidrule(lr){2-4} \cmidrule(lr){5-8}  \cmidrule(lr){9-11} \cmidrule(lr){12-16} \cmidrule(lr){17-20}
    & Gauss. & Shot & Impul. & Defoc. & Glass & Motion & Zoom & Snow & Frost & Fog & Brit. & Contr. & Elastic & Pixel & JPEG & Acc (↑) & AUROC (↑) & ECE (↓) &  MCE (↓)  \\
    \midrule
    Source & 44.99 & 46.80 & 47.25 & 82.05 & 44.04 & 76.54 & 84.05 & 78.80 & 75.59 & 67.81 & \textbf{87.67} & 63.71 & 65.82 & 55.31 & 64.98 & 65.69 & 80.79 & 16.16 & 37.95 \\
    TENT\cite{ftent} & 51.29 & 46.23 & 35.02 & 79.23 & 6.38 & 39.17 & 27.15 & 2.90 & 1.41 & 1.07 & 4.34 & 1.18 & 1.00 & 1.01 & 1.02 & 19.89 & 43.05 & 76.34 & 80.71 \\
    SAR\cite{sar} & 59.50 & 68.16 & 65.53 & 81.27 & 62.48 & 77.42 & 84.20 & 78.73 & 80.69 & 72.69 & 85.78 & 76.69 & 70.71 & 76.14 & 66.01 & 73.73 & 47.42 & 16.89 & 39.24 \\
    TEA\cite{tea} & 49.02 & 56.23 & 56.46 & \textbf{82.57} & 54.46 & 77.71 & 84.76 & 79.24 & 79.66 & 72.66 & 86.80 & 73.37 & 69.61 & 72.98 & 64.71 & 70.68 & 83.98 & 17.37 & 37.38 \\
    TENT(APMS) & \textbf{61.17} & \textbf{70.30} & \textbf{69.19} & 81.53 & \textbf{66.35} & \textbf{78.47} & \textbf{84.87} & \textbf{80.47} & \textbf{82.66} & \textbf{75.43} & 86.46 & \textbf{78.25} & \textbf{71.89} & \textbf{78.27} & \textbf{67.22} & \textbf{75.50} & \textbf{85.81} & 17.68 & 43.53 \\
    SAR(APMS) & 51.91 & 60.83 & 56.97 & 81.59 & 56.79 & 76.58 & 83.64 & 78.99 & 79.75 & 71.00 & 87.02 & 68.75 & 69.17 & 73.01 & 65.89 & 70.79 & 66.02 & 16.35 & 36.89 \\
    TEA(APMS) & 48.67 & 56.58 & 56.05 & 82.16 & 56.79 & 76.94 & 84.00 & 79.03 & 79.96 & 72.97 & 86.18 & 73.35 & 68.87 & 72.24 & 62.78 & 70.44 & 83.65 & \textbf{15.70} & \textbf{36.60} \\
    \end{tabular}
    \end{latin}
    \end{adjustbox}
\end{table}

\begin{table}[H]
    \centering
    \caption[مقایسه‌ی میانگین عملکرد در مجموعه داده‌ی \lr{CIFAR100-C} در حالت مستمر]{مقایسه‌ی میانگین عملکرد بر روی مدل ViT در مجموعه داده‌ی \lr{CIFAR100-C} در حالت مستمر}
    \label{table:t8}
    \begin{adjustbox}{width=1.5\textwidth}
    \begin{latin}
    \begin{tabular}{lccccccccccccccc|cccc}
    \toprule
    Time & \multicolumn{15}{l|}{$t \xrightarrow{\hspace*{23cm}}$} & \\ 
    \midrule
    \multirow{2}{*}{Method} & \multicolumn{3}{c}{Noise} & \multicolumn{4}{c}{Blur} & \multicolumn{3}{c}{Weather} & \multicolumn{5}{c}{Digital} & \multicolumn{4}{c}{Avg} \\
    \cmidrule(lr){2-4} \cmidrule(lr){5-8}  \cmidrule(lr){9-11} \cmidrule(lr){12-16} \cmidrule(lr){17-20}
    & Gauss. & Shot & Impul. & Defoc. & Glass & Motion & Zoom & Snow & Frost & Fog & Brit. & Contr. & Elastic & Pixel & JPEG & Acc (↑) & AUROC (↑) & ECE (↓) &  MCE (↓)  \\
    \midrule
    Source & 44.99 & 46.80 & 47.25 & \textbf{82.05} & 44.04 & \textbf{76.54} & \textbf{84.05} & \textbf{78.80} & \textbf{75.59} & \textbf{67.81} & \textbf{87.67} & \textbf{63.71} & \textbf{65.82} & \textbf{55.31} & \textbf{64.98} & \textbf{65.69} & \textbf{80.79} & \textbf{16.16} & 37.95 \\
    TENT\cite{ftent} & 36.54 & 12.74 & 9.09 & 22.56 & 2.08 & 8.66 & 6.27 & 1.38 & 1.08 & 1.01 & 1.76 & 1.08 & 1.00 & 1.02 & 1.01 & 7.15 & 42.73 & 91.51 & 92.93 \\
    SAR\cite{sar} & 37.27 & 32.80 & 25.49 & 34.00 & 22.00 & 29.66 & 32.26 & 30.05 & 30.55 & 27.57 & 33.14 & 28.36 & 26.79 & 28.67 & 25.11 & 29.58 & 48.76 & 65.76 & 74.50 \\
    TEA\cite{tea} & 42.19 & 36.17 & 20.17 & 25.73 & 11.69 & 16.41 & 17.76 & 16.65 & 16.73 & 15.33 & 18.16 & 15.47 & 14.72 & 15.40 & 13.74 & 19.75 & 54.59 & 76.84 & 82.16 \\
    TENT(APMS) & \textbf{63.63} & \textbf{62.20} & \textbf{55.34} & 65.04 & \textbf{52.55} & 61.97 & 66.40 & 62.03 & 63.82 & 50.03 & 66.03 & 49.58 & 41.19 & 31.91 & 27.38 & 54.61 & 75.04 & 40.71 & 58.69 \\
    SAR(APMS) & 50.70 & 50.70 & 48.62 & 71.67 & 47.49 & 67.40 & 73.54 & 69.21 & 69.01 & 62.06 & 76.33 & 60.07 & 59.88 & 51.95 & 49.67 & 60.55 & 51.64 & 16.26 & \textbf{34.91} \\
    TEA(APMS) & 54.24 & 49.60 & 38.89 & 49.59 & 38.81 & 47.36 & 50.90 & 48.21 & 49.21 & 45.46 & 51.73 & 47.16 & 41.21 & 31.38 & 20.25 & 44.27 & 70.33 & 49.98 & 64.60 \\
    \end{tabular}
    \end{latin}
    \end{adjustbox}
\end{table}
\end{landscape}

\begin{landscape}
\begin{table}[H]
    \centering
    \caption[نتیجه‌ی اعمال الگوریتم شناسایی لایه‌ی بهینه در مجموعه داده‌ی \lr{Cifar100}.]{نتایج اعمال الگوریتم شناسایی لایه‌ی بهینه بر روی مجموعه‌داده‌ی \lr{CIFAR100} و مدل مبتنی بر مبدل بینایی. در این مدل، لایه‌ی ۰ متناظر با خروجی لایه‌ی \lr{Patch Embedding} و دوازده لایه‌ی بعدی متناظر با خروجی‌ لایه‌های دوازده \lr{Transformer Encoder} هستند. بر اساس مقادیر امتیاز جرم تقریبی، لایه‌ی چهارم به‌عنوان لایه‌ی بهینه انتخاب می‌شود.}    
    \label{table:alg2_results}
    \begin{latin}
    \begin{tabular}{c c c c c c c c c c c c c c}
    \toprule
    Layer Index & 0 & 1 & 2 & 3 & 4 & 5 & 6 & 7 & 8 & 9 & 10 & 11 & 12 \\
    \midrule
    APMS        & 0.0117 & 0.0148 & 0.0190 & 0.0406 & 0.0508 & 0.0496 & 0.0440 & 0.0368 & 0.0246 & 0.0196 & 0.0039 & 0.0024 & 0.0272 \\
    \bottomrule
    \end{tabular}
    \end{latin}
\end{table}

\begin{table}[H]
    \centering
    \caption[مقایسه‌ی عملکرد در مجموعه داده‌ی \lr{CIFAR100-C} با شناسایی لایه‌ی بهینه.]{مقایسه‌ی عملکرد بر روی مدل ViT در مجموعه داده‌ی \lr{CIFAR100-C} با شناسایی لایه‌ی بهینه‌ی حاصل از الگوریتم۲.}
    \label{table:t9}
    \begin{adjustbox}{width=1.5\textwidth}
    \begin{latin}
    \begin{tabular}{lccccccccccccccc|cccc}
    \toprule
    \multirow{2}{*}{Method} & \multicolumn{3}{c}{Noise} & \multicolumn{4}{c}{Blur} & \multicolumn{3}{c}{Weather} & \multicolumn{5}{c}{Digital} & \multicolumn{4}{c}{Avg} \\
    \cmidrule(lr){2-4} \cmidrule(lr){5-8}  \cmidrule(lr){9-11} \cmidrule(lr){12-16} \cmidrule(lr){17-20}
    & Gauss. & Shot & Impul. & Defoc. & Glass & Motion & Zoom & Snow & Frost & Fog & Brit. & Contr. & Elastic & Pixel & JPEG & Acc (↑) & AUROC (↑) & ECE (↓) &  MCE (↓)  \\
    \midrule
    Source & 44.99 & 46.80 & 47.25 & 82.05 & 44.04 & 76.54 & 84.05 & 78.80 & 75.59 & 67.81 & 87.67 & 63.71 & 65.82 & 55.31 & 64.98 & 65.69 & 80.79 & 16.16 & \textbf{37.95} \\
    TENT(APMS: Layer0) & \textbf{66.34} & \textbf{69.57} & \textbf{71.97} & \textbf{84.72} & \textbf{71.26} & \textbf{83.68} & \textbf{86.75} & \textbf{83.38} & \textbf{83.46} & \textbf{79.89} & \textbf{88.93} & \textbf{85.50} & \textbf{77.99} & \textbf{83.96} & \textbf{73.48} & \textbf{79.39} & \textbf{87.01} & \textbf{13.80} & 38.53 \\
    TENT(APMS: Layer4) & 65.29 & 19.93 & 70.88 & 84.62 & 68.54 & 82.55 & 86.29 & 82.74 & 82.20 & 78.00 & 88.92 & 84.50 & 77.09 & 83.07 & 72.78 & 75.16 & 83.48 & 18.40 & 42.79 \\
    \end{tabular}
    \end{latin}
    \end{adjustbox}
\end{table}
\end{landscape}

\زیرقسمت{تحلیل پیچیدگی}

یکی از جنبه‌های مهم در ارزیابی روش‌های انطباق زمان آزمون، بررسی پیچیدگی محاسباتی آن‌هاست. پیچیدگی هر روش را می‌توان بر اساس تعداد {جلورو\پانویس{Forward-pass}} و {عقب‌رو\پانویس{Backward-pass}} موردنیاز در هر گام انطباق تحلیل کرد.  

در گام زمانی $t$ دسته‌ی ورودی به صورت $X_t = \{x_1, \dots, x_m\}$ به مدل برای پیش‌بینی و انطباق داده می‌شود.  

در روش \lr{TENT}، ابتدا $m$ بار جلورو روی $X_t$ انجام می‌شود تا آنتروپی خروجی‌ها محاسبه شود و سپس برای به‌روزرسانی وزن‌ها بر اساس گرادیان آنتروپی $m$ عقب‌رو انجام می‌شود؛ بنابراین تعداد جلورو نهایی برابر $m$ و تعداد عقب‌رو نهایی برابر $m$ خواهد بود.  

در نسخه‌ی بهبود یافته‌ی روش \lr{TENT} نیز ابتدا $m$ بار جلورو روی $X_t$ انجام می‌شود، سپس یک مرحله پالایش نمونه‌ها صورت می‌گیرد و زیرمجموعه‌ی $X'_t = \{ x \in X_t \mid S_{\mathrm{ent}}(x) = 1\}$ به‌دست می‌آید؛ پس از آن ابتدا برای محاسبه وزن‌های $W(x)$ یک عقب‌رو و سپس برای به‌روزرسانی پارامترهای شبکه یک عقب‌رو دیگر روی مجموعه $X'_t$ انجام می‌شود، بنابراین تعداد جلورو نهایی برابر $m$ و تعداد عقب‌رو نهایی برابر $2 |X'_t|$ خواهد بود.  

در روش \lr{SAR} ابتدا $m$ بار جلورو روی $X_t$ انجام می‌شود، سپس زیرمجموعه‌ی $X'_t = \{ x \in X_t \mid S_{\mathrm{ent}}(x) = 1\}$ به‌دست می‌آید و مرحله‌ی بیشینه‌سازی رابطه‌ی \ref{eq:sar} روی آن انجام می‌شود که شامل $|X'_t|$ عقب‌رو است؛ پس از به‌روزرسانی، پالایش دوم انجام شده و مجموعه‌ی $X''_t = \{ x \in X'_t \mid S_{\mathrm{ent}}(x) = 1\}$ به دست می‌آید و روی آن یک عقب‌رو دیگر انجام می‌شود؛ بنابراین تعداد جلورو نهایی برابر $m + |X'_t|$ و تعداد عقب‌رو نهایی برابر $|X'_t| + |X''_t|$ خواهد بود.  

در نسخه‌ی بهبود یافته‌ی روش \lr{SAR}، تنها یک عقب‌رو بیشتر برای محاسبه‌ی وزن‌های $W(x)$ بعد از پالایش اولیه نیاز است. بنابراین تعداد جلورو نهایی برابر $m + |X'_t|$ و تعداد عقب‌رو نهایی برابر $|X'_t| + |X''_t| + |X_t|$ خواهد بود.  

در روش \lr{TEA} ابتدا $m$ بار جلورو روی $X_t$ انجام می‌شود، سپس به ازای هر نمونه با $k$ گام \lr{SGLD} نمونه‌های مصنوعی ایجاد می‌شوند که شامل $k \cdot m$ بار جلورو و $k \cdot m$ بار عقب‌رو برای محاسبه گرادیان انرژی است، و در نهایت یک عقب‌رو برای حل تابع هدف اصلی روی نمونه‌های واقعی و مصنوعی انجام می‌شود؛ بنابراین تعداد جلورو نهایی برابر $(k+1) m$ و تعداد عقب‌رو نهایی برابر $(k+2) m$ خواهد بود.  

در نسخه‌ی بهبود یافته‌ی روش \lr{TEA} ابتدا $m$ بار جلورو روی دسته‌ی ورودی $X_t$ انجام می‌شود، سپس یک مرحله پالایش نمونه‌ها صورت می‌گیرد و زیرمجموعه‌ی $X'_t$ به‌دست می‌آید. سپس به ازای هر نمونه در $X'_t$ با $k$ گام \lr{SGLD} نمونه‌های مصنوعی ایجاد می‌شود که شامل $k \cdot |X'_t|$ بار جلورو و $k \cdot |X'_t|$ بار عقب‌رو برای محاسبه گرادیان انرژی است؛ در نهایت یک بار عقب‌رو برای حل تابع هدف اصلی روی نمونه‌های زمان آزمون و نمونه‌های مصنوعی ایجادشده از آن‌ها انجام می‌شود، بنابراین تعداد جلورو نهایی برابر $m + k|X'_t|$ و تعداد عقب‌رو نهایی برابر $2 |X'_t| + k|X'_t| + 1 = (k+2)|X'_t| + 1$ خواهد بود. جمع‌بندی کلی تحلیل پیچیدگی در جدول \ref{table:complexity} آورده شده‌است.

\begin{table}[t]
    \centering
    \small
    \caption[تحلیل پیچیدگی]{تحلیل پیچیدگی روش‌های  پایه‌ای و روش پیشنهادی}
    \label{table:complexity}
    \setlength{\tabcolsep}{6pt}
    \begin{tabular}{p{4cm} p{3.5cm} p{3.5cm}}
    \toprule
    \textbf{روش} & \textbf{\RL{\#جلورو}} & \textbf{\RL{\#عقب‌رو}}\\ \midrule
    \lr{TENT}\cite{ftent}         & $m$              & $m$ \\
    \lr{SAR}\cite{sar}          & $m + |X'_t|$     & $|X'_t| + |X''_t|$ \\
    \lr{TEA}\cite{tea}         & $(k+1)m$         & $(k+2)m$ \\
    \lr{TENT (APMS)}  & $m$              & $2 |X'_t|$ \\
    \lr{SAR (APMS)}   & $m + |X'_t|$     & $|X'_t| + |X''_t| + |X_t|$ \\
    \lr{TEA (APMS)}   & $m + k|X'_t|$    & $(k+2)|X'_t| + 1$ \\
    \bottomrule
    \end{tabular}
\end{table}


\زیرقسمت{تحلیل پایداری و حساسیت}

در این بخش، به بررسی پایداری و حساسیت روش‌های پیشنهادی و پایه‌ای پرداخته می‌شود. هدف اصلی تحلیل، سنجش میزان ثبات عملکرد مدل در مواجهه با تغییرات مختلف داده‌ها و همچنین بررسی حساسیت آن نسبت به انتخاب پارامترها و شرایط متفاوت آزمایش است. این تحلیل، مکمل نتایج کمی جداول عملکرد است و دید دقیق‌تری نسبت به قابلیت اعتماد و قابلیت تعمیم روش‌ها فراهم می‌کند.

برای این منظور، روش‌های پایه‌ای \lr{TENT}، \lr{SAR} و \lr{TEA} در نظر گرفته شده و عملکرد نسخه‌های بهبود یافته با تزریق معیار امتیاز جرم تقریبی نیز بررسی می‌شود. تحلیل‌ها شامل مقایسه‌ی تغییرات نرخ خطا و معیارهای مورد استفاده در شرایط مختلف است تا نشان داده شود که تزریق معیار پیشنهادی چه تأثیری بر کاهش واریانس عملکرد و افزایش ثبات مدل دارد.

نمودارهای ارائه شده در ادامه، این تحلیل‌ها را برای مجموعه داده‌های \lr{CIFAR10-C} و \lr{CIFAR100-C} نشان می‌دهند و به‌طور مستقیم مقایسه‌ی پایداری روش‌های پایه‌ای و نسخه‌های بهبود یافته با معیار پیشنهادی را ممکن می‌سازند. این نمودارها نشان می‌دهند که چگونه معیار امتیاز جرم تقریبی، نه تنها دقت مدل را ارتقا می‌دهد، بلکه حساسیت آن به تغییرات داده‌ها و پارامترها را نیز کاهش می‌دهد و موجب استحکام فرایند انطباق می‌شود.

نمودارهای ارائه‌شده در شکل  \ref{fig:robustness_all} به بررسی استحکام و پایداری روش‌ها نسبت به نرخ‌های مختلف یادگیری اختصاص دارند. در هر نمودار، عملکرد روش پایه با نسخه‌ی بهبود یافته‌ی آن که از معیار امتیاز جرم تقریبی بهره می‌برد، مقایسه شده است. محور افقی نرخ یادگیری و محور عمودی نرخ خطا را نشان می‌دهد.
این مقایسه‌ی مستقیم امکان ارزیابی ثبات عملکرد روش‌ها را در مواجهه با تغییرات مهم در فراپارامترها فراهم می‌کند. همان‌طور که مشاهده می‌شود، نسخه‌های بهبود یافته نه تنها نرخ خطای پایین‌تری ارائه می‌دهند، بلکه تغییرات عملکرد آن‌ها در پاسخ به نرخ‌های یادگیری متفاوت نیز کمتر است؛ به عبارت دیگر، روش پیشنهادی حساسیت کمتری به انتخاب دقیق نرخ یادگیری دارد و پایداری بالاتری را به نمایش می‌گذارد.\\
هم‌چنین نمودار‌های ارائه‌شده در شکل \ref{fig:sensivity_all} به بررسی حساسیت روش‌ها نسبت به تغییرات جزئی در نرخ یادگیری اختصاص دارند. در این تحلیل، عملکرد روش پایه با نسخه‌ی بهبود یافته‌ی آن که از معیار امتیاز جرم تقریبی بهره می‌برد، مقایسه شده است. این بررسی امکان ارزیابی میزان حساسیت روش‌ها به تغییرات کوچک در فراپارامترها را فراهم می‌کند. همان‌طور که مشاهده می‌شود، نسخه‌های بهبود یافته نه تنها نرخ خطای پایین‌تری ارائه می‌کنند، بلکه تغییرات عملکرد آن‌ها در مواجهه با تغییرات جزئی نرخ یادگیری نیز کمتر است؛ به عبارت دیگر، روش پیشنهادی نسبت به نوسانات کوچک در پارامترها مقاوم‌تر بوده و ثبات بیشتری در عملکرد خود نشان می‌دهد.

به طور کلی، این تحلیل‌ها نشان می‌دهد که تزریق معیار امتیاز جرم تقریبی، علاوه بر ارتقای دقت، هم حساسیت روش‌ها نسبت به تغییرات کوچک پارامترها و هم آسیب‌پذیری آن‌ها در مقابل تغییرات مهم فراپارامترها را کاهش داده و قابلیت اعتماد مدل را به‌طور قابل توجهی افزایش می‌دهد.
\begin{figure}[t]
	\centering
	\subfigure[]{
		\includegraphics[width=0.48\linewidth]{images/robustness/cifar100c/TENT.png}
	}
	\hfill
	\subfigure[]{
		\includegraphics[width=0.48\linewidth]{images/robustness/cifar100c/TENT-continual.png}
	}
	
	\subfigure[]{
		\includegraphics[width=0.48\linewidth]{images/robustness/cifar100c/sar.png}
	}
	\hfill
	\subfigure[]{
		\includegraphics[width=0.48\linewidth]{images/robustness/cifar100c/sar-continual.png}
	}
	
	\subfigure[]{
		\includegraphics[width=0.48\linewidth]{images/robustness/cifar100c/tea.png}
	}
	\hfill
	\subfigure[]{
		\includegraphics[width=0.48\linewidth]{images/robustness/cifar100c/tea-continual.png}
	}
	
	\caption[تحلیل پایداری]{بررسی پایداری روش‌های پایه‌ای و نسخه‌های بهبود یافته با تزریق امتیاز جرم تقریبی نسبت به نرخ‌های مختلف یادگیری؛ ستون راست: انطباق استاندارد، ستون چپ: انطباق مستمر. هر سطر مربوط به یک روش است (\lr{TENT}، \lr{SAR}، \lr{TEA}).}
	\label{fig:robustness_all}
\end{figure}

\begin{figure}[t]
	\centering
	\subfigure[]{
		\includegraphics[width=0.48\linewidth]{images/sensivity/cifar100c/TENT1.png}
	}
	\hfill
	\subfigure[]{
		\includegraphics[width=0.48\linewidth]{images/sensivity/cifar100c/TENT-continual.png}
	}
	
	\subfigure[]{
		\includegraphics[width=0.48\linewidth]{images/sensivity/cifar100c/sar.png}
	}
	\hfill
	\subfigure[]{
		\includegraphics[width=0.48\linewidth]{images/sensivity/cifar100c/sar-continual.png}
	}
	
	\subfigure[]{
		\includegraphics[width=0.48\linewidth]{images/sensivity/cifar100c/tea.png}
	}
	\hfill
	\subfigure[]{
		\includegraphics[width=0.48\linewidth]{images/sensivity/cifar100c/tea-continual.png}
	}
	\caption[تحلیل حساسیت]{بررسی حساسیت روش‌های پایه‌ای و نسخه‌های بهبود یافته با تزریق امتیاز جرم تقریبی نسبت به نرخ‌های مختلف یادگیری؛ ستون راست: انطباق استاندارد، ستون چپ: انطباق مستمر. هر سطر مربوط به یک روش است (\lr{TENT}، \lr{SAR}، \lr{TEA}).}
	\label{fig:sensivity_all}
\end{figure}
\FloatBarrier

\زیرقسمت{بصری‌سازی Grad-CAM}

برای تحلیل بصری تأثیر فرآیند انطباق بر تصمیمات مدل، از روش \lr{Grad-CAM} استفاده شده است. این ابزار با برجسته‌سازی نواحی از تصویر که بیشترین نقش را در پیش‌بینی دارند، امکان مشاهده‌ی مستقیم تمرکز مدل را فراهم می‌سازد و کمک می‌کند تا درک شهودی‌تری از نحوه‌ی استخراج ویژگی‌ها و تغییرات ایجادشده در اثر انطباق به دست آید. چنین تحلیلی به‌ویژه در شرایطی که داده‌ها دچار اغتشاش یا تخریب شده‌اند، ارزشمند است زیرا نشان می‌دهد مدل در عمل بر چه بخش‌هایی از تصویر تکیه می‌کند.

بررسی نتایج شکل \ref{fig:gradcam_visualization}  نشان می‌دهد که در مدل منبع و در روش پایه‌ای \lr{TENT} که انطباق بدون هیچ‌گونه کنترل یا قید اضافی باعث می‌شود نقشه‌های توجه گرایش به تمرکز بر بخش‌های نامرتبط داشته باشند. این موضوع عملاً منجر به استخراج ویژگی‌های اشتباه شده و در ادامه با ایجاد انباشت خطا، خطر فروپاشی مدل نیز افزایش می‌یابد. چنین رفتاری نشان می‌دهد که هرچند انطباق می‌تواند بهبودهایی به همراه داشته باشد، در غیاب سازوکارهای کنترلی احتمال انحراف مدل از مسیر درست بسیار بالاست.

در مقابل، نسخه‌ی بهبود یافته که در آن معیار امتیاز جرم تقریبی تزریق شده است، رفتار پایدارتری نشان می‌دهد. نقشه‌های حاصل از \lr{Grad-CAM} در این حالت بیانگر آن است که تمرکز مدل به‌جای بخش‌های پراکنده یا نامرتبط، به سمت نواحی معنادار و مرتبط با هدف هدایت می‌شود. این مسئله نه‌تنها نشان‌دهنده‌ی استخراج صحیح‌تر ویژگی‌هاست، بلکه به‌طور غیرمستقیم تأیید می‌کند که مدل توانایی مقابله با خطاهای انباشته و جلوگیری از فروپاشی را دارد.

در مجموع، ترکیب این تحلیل کیفی با ارزیابی‌های کمی، تصویری کامل از عملکرد روش پیشنهادی ارائه می‌دهد. نتایج نشان می‌دهند که رویکرد بهبود یافته علاوه بر افزایش دقت، ثبات و پایداری بیشتری را نیز تضمین می‌کند و شهودی روشن از چگونگی کنترل فرآیند انطباق و کاهش خطر فروپاشی مدل در حالتهای انطباق برخط زمان آزمون در اختیار قرار می‌دهد.

\begin{figure}[t]
	\centering
	\subfigure[]{
		\includegraphics[width=0.9\linewidth]{images/gradcam.png}
	}
	\caption[بصری‌سازی \lr{Grad-CAM}]{بصری‌سازی Grad-Cam: برای هر تصویر ورودی (ستون اول از سمت چپ)، ستون میانی نقشه‌ی توجهی مدل پایه‌ای \lr{TENT} و ستون سمت راست نقشه‌ی توجهی نسخه‌ی بهبود یافته \lr{TENT} با تزریق معیار امتیاز جرم تقریبی را نشان می‌دهد.}
	\label{fig:gradcam_visualization}
\end{figure}
\FloatBarrier

\قسمت{جمع‌بندی}

در این فصل، عملکرد روش پیشنهادی در حالتهای انطباق زمان آزمون و انطباق مستمر زمان آزمون به‌طور جامع بررسی شد. نتایج حاصل از آزمایش‌ها نشان داد که تزریق معیار امتیاز جرم تقریبی به روش‌های پایه‌ای، علاوه بر ارتقای دقت پیش‌بینی‌ها، نقش مهمی در افزایش پایداری و کاهش حساسیت روش‌ها نسبت به تغییرات فراپارامترها دارد. تحلیل‌های کمی و نمودارهای پایداری و حساسیت نشان دادند که نسخه‌های بهبود یافته، عملکردی همگن‌تر و با واریانس کمتر ارائه می‌دهند و از فروپاشی مدل در طول به‌روزرسانی‌های مکرر جلوگیری می‌کنند.

تحلیل بصری با استفاده از \lr{Grad-CAM} نیز تأیید کرد که انطباق با معیار پیشنهادی، تمرکز مدل را به سمت نواحی معنادار تصویر هدایت می‌کند و از اتلاف تمرکز بر بخش‌های نامرتبط جلوگیری می‌کند. این رفتار مدل نه تنها منجر به استخراج ویژگی‌های صحیح‌تر می‌شود، بلکه توانایی مقابله با انباشت خطا و حفظ ساختار اصلی شبکه را نیز تضمین می‌کند.

به طور کلی، نتایج این فصل نشان می‌دهد که روش پیشنهادی قادر است در شرایط متغیر و چالش‌برانگیز داده‌ها، عملکردی دقیق، پایدار و مقاوم ارائه دهد. این یافته‌ها اهمیت بالای انتخاب معیار معتبر در فرآیند انطباق زمان آزمون و تأثیر آن بر افزایش قابلیت اعتماد مدل را برجسته می‌سازد و پایه محکمی برای کاربردهای عملی روش در محیط‌های پویا و داده‌های غیرایستا فراهم می‌آورد.